\section{BAB I}\label{bab-i}

\section{PENDAHULUAN}\label{pendahuluan}

\subsection{1.1 Latar Belakang}\label{latar-belakang}

Perkembangan teknologi informasi telah memberikan banyak kemudahan dalam berbagai bidang usaha, termasuk pada sektor makanan dan minuman (Food and Beverage/F\&B). Salah satu kebutuhan penting dalam operasional kafe dan restoran adalah sistem yang mampu membantu proses transaksi penjualan secara cepat, akurat, dan terintegrasi.

Pada banyak usaha kecil hingga menengah, proses pencatatan transaksi masih dilakukan secara manual atau menggunakan aplikasi yang memiliki keterbatasan fitur. Hal tersebut dapat menyebabkan kesalahan pencatatan, kesulitan dalam pengelolaan stok, serta keterlambatan dalam penyusunan laporan penjualan. Selain itu, data transaksi yang tersimpan sering kali hanya digunakan sebagai arsip tanpa dimanfaatkan untuk menghasilkan informasi yang dapat mendukung pengambilan keputusan bisnis.

Untuk mengatasi permasalahan tersebut, dikembangkan sebuah sistem Point of Sale (POS) berbasis web bernama \textbf{Brew \& Bytes POS}. Sistem ini dirancang untuk membantu proses transaksi penjualan, pengelolaan produk, manajemen stok, pengelolaan pengguna, serta penyusunan laporan penjualan secara terintegrasi.

Selain menyediakan fungsi kasir, sistem ini juga dilengkapi dengan fitur pembayaran digital menggunakan QRIS melalui integrasi Payment Gateway Midtrans. Sistem juga menerapkan metode FIFO (First In First Out) dalam pengelolaan stok agar pergerakan persediaan dapat tercatat dengan lebih baik.

Sebagai nilai tambah, sistem ini mengimplementasikan algoritma Apriori untuk melakukan analisis pola pembelian pelanggan (Market Basket Analysis). Hasil analisis tersebut dapat digunakan untuk mengetahui produk yang sering dibeli secara bersamaan sehingga dapat membantu pemilik usaha dalam menyusun strategi promosi dan paket bundling produk.

Berdasarkan kebutuhan tersebut, dibuatlah proyek \textbf{Brew \& Bytes POS}, yaitu sistem Point of Sale berbasis web yang terintegrasi dengan fitur manajemen inventori, pembayaran digital, dan analisis pola pembelian pelanggan.

\begin{center}\rule{0.5\linewidth}{0.5pt}\end{center}

\subsection{1.2 Tujuan Proyek}\label{tujuan-proyek}

Tujuan dari pengembangan proyek ini adalah sebagai berikut:

\begin{enumerate}
\def\labelenumi{\arabic{enumi}.}
\tightlist
\item
  Membangun sistem Point of Sale berbasis web untuk membantu proses transaksi penjualan.
\item
  Mempermudah pengelolaan data produk dan kategori.
\item
  Mengelola stok barang secara terintegrasi menggunakan metode FIFO.
\item
  Menyediakan laporan penjualan yang dapat diakses secara mudah.
\item
  Mengintegrasikan pembayaran digital menggunakan QRIS.
\item
  Menerapkan algoritma Apriori untuk menganalisis pola pembelian pelanggan.
\end{enumerate}

\begin{center}\rule{0.5\linewidth}{0.5pt}\end{center}

\subsection{1.3 Ruang Lingkup Proyek}\label{ruang-lingkup-proyek}

Ruang lingkup proyek yang dikembangkan meliputi:

\begin{enumerate}
\def\labelenumi{\arabic{enumi}.}
\tightlist
\item
  Sistem berbasis web yang dapat diakses melalui browser.
\item
  Manajemen pengguna dengan hak akses Admin, Manager, dan Cashier.
\item
  Manajemen produk dan kategori.
\item
  Transaksi penjualan melalui modul Point of Sale (POS).
\item
  Pembayaran menggunakan metode tunai dan QRIS.
\item
  Pengelolaan stok barang menggunakan metode FIFO.
\item
  Laporan penjualan dan statistik bisnis.
\item
  Analisis pola pembelian pelanggan menggunakan algoritma Apriori.
\item
  Dukungan mode offline untuk transaksi tertentu.
\end{enumerate}

\begin{center}\rule{0.5\linewidth}{0.5pt}\end{center}

\subsection{1.4 Manfaat Proyek}\label{manfaat-proyek}

\subsubsection{Bagi Pengguna}\label{bagi-pengguna}

\begin{enumerate}
\def\labelenumi{\arabic{enumi}.}
\tightlist
\item
  Mempermudah proses transaksi penjualan.
\item
  Mengurangi kesalahan pencatatan transaksi.
\item
  Mempermudah pengelolaan stok dan produk.
\item
  Mempercepat proses penyusunan laporan.
\end{enumerate}

\subsubsection{Bagi Pemilik Usaha}\label{bagi-pemilik-usaha}

\begin{enumerate}
\def\labelenumi{\arabic{enumi}.}
\tightlist
\item
  Memperoleh informasi penjualan secara real-time.
\item
  Memantau kondisi stok dengan lebih mudah.
\item
  Memanfaatkan analisis pola pembelian pelanggan untuk strategi pemasaran.
\item
  Mendukung pengambilan keputusan berdasarkan data transaksi.
\end{enumerate}

\subsubsection{Bagi Pengembang}\label{bagi-pengembang}

\begin{enumerate}
\def\labelenumi{\arabic{enumi}.}
\tightlist
\item
  Menerapkan teknologi web modern dalam pengembangan aplikasi.
\item
  Mengimplementasikan integrasi payment gateway.
\item
  Mengimplementasikan algoritma Apriori pada studi kasus nyata.
\item
  Menambah pengalaman dalam pengembangan sistem informasi berbasis web.
\end{enumerate}

\begin{center}\rule{0.5\linewidth}{0.5pt}\end{center}

\subsection{1.5 Gambaran Umum Sistem}\label{gambaran-umum-sistem}

Brew \& Bytes POS merupakan aplikasi Point of Sale berbasis web yang dibangun menggunakan Laravel, React, dan MySQL. Sistem ini dirancang untuk membantu operasional kafe atau restoran melalui pengelolaan transaksi, inventori, pembayaran digital, serta analisis data penjualan.

Fitur utama yang tersedia dalam sistem antara lain:

\begin{itemize}
\tightlist
\item
  Login dan manajemen pengguna
\item
  Dashboard bisnis
\item
  Point of Sale (POS)
\item
  Manajemen produk
\item
  Manajemen kategori
\item
  Pembayaran QRIS
\item
  Laporan penjualan
\item
  Market Basket Analysis
\item
  Sinkronisasi offline
\end{itemize}

Dengan adanya sistem ini diharapkan proses operasional bisnis dapat berjalan lebih efektif, efisien, dan terintegrasi dalam satu platform.

\section{BAB II}\label{bab-ii}

\section{GAMBARAN UMUM DAN PERANCANGAN SISTEM}\label{gambaran-umum-dan-perancangan-sistem}

\subsection{2.1 Gambaran Umum Sistem}\label{gambaran-umum-sistem-1}

Brew \& Bytes POS merupakan sistem Point of Sale (POS) berbasis web yang dirancang untuk membantu operasional usaha kafe dan restoran. Sistem ini menyediakan berbagai fitur yang mendukung proses bisnis mulai dari pengelolaan produk, transaksi penjualan, pembayaran digital, pengelolaan stok, hingga analisis pola pembelian pelanggan.

Sistem dibangun menggunakan framework Laravel sebagai backend, React sebagai frontend, serta MySQL sebagai basis data. Selain itu, sistem terintegrasi dengan layanan Midtrans untuk mendukung pembayaran QRIS dan menerapkan algoritma Apriori untuk menghasilkan analisis Market Basket Analysis.

Tujuan utama sistem adalah meningkatkan efisiensi operasional bisnis serta membantu pengambilan keputusan berdasarkan data transaksi yang tersimpan.

\begin{center}\rule{0.5\linewidth}{0.5pt}\end{center}

\subsection{2.2 Arsitektur Sistem}\label{arsitektur-sistem}

Brew \& Bytes POS menggunakan arsitektur client-server.

Pada arsitektur ini, pengguna mengakses aplikasi melalui browser yang berkomunikasi dengan server aplikasi. Server bertanggung jawab untuk memproses data, menjalankan logika bisnis, berinteraksi dengan database, serta mengelola integrasi dengan layanan pihak ketiga.

\subsubsection{Diagram Arsitektur Sistem}\label{diagram-arsitektur-sistem}

\begin{Shaded}
\begin{Highlighting}[]
\NormalTok{+{-}{-}{-}{-}{-}{-}{-}{-}{-}{-}{-}{-}{-}{-}{-}{-}{-}{-}+}
\NormalTok{|      User        |}
\NormalTok{+{-}{-}{-}{-}{-}{-}{-}{-}{-}{-}{-}{-}{-}{-}{-}{-}{-}{-}+}
\NormalTok{         │}
\NormalTok{         ▼}
\NormalTok{+{-}{-}{-}{-}{-}{-}{-}{-}{-}{-}{-}{-}{-}{-}{-}{-}{-}{-}+}
\NormalTok{| React + Inertia  |}
\NormalTok{|    Frontend      |}
\NormalTok{+{-}{-}{-}{-}{-}{-}{-}{-}{-}{-}{-}{-}{-}{-}{-}{-}{-}{-}+}
\NormalTok{         │}
\NormalTok{         ▼}
\NormalTok{+{-}{-}{-}{-}{-}{-}{-}{-}{-}{-}{-}{-}{-}{-}{-}{-}{-}{-}+}
\NormalTok{| Laravel Backend  |}
\NormalTok{+{-}{-}{-}{-}{-}{-}{-}{-}{-}{-}{-}{-}{-}{-}{-}{-}{-}{-}+}
\NormalTok{         │}
\NormalTok{ ┌───────┼─────────┐}
\NormalTok{ ▼       ▼         ▼}
\NormalTok{MySQL  Midtrans  Analytics}
\NormalTok{Database Payment  Apriori}
\end{Highlighting}
\end{Shaded}

\subsubsection{Penjelasan Arsitektur}\label{penjelasan-arsitektur}

\begin{itemize}
\tightlist
\item
  React digunakan untuk membangun antarmuka pengguna.
\item
  Inertia.js digunakan sebagai penghubung antara frontend dan backend.
\item
  Laravel menangani logika bisnis aplikasi.
\item
  MySQL digunakan untuk menyimpan seluruh data sistem.
\item
  Midtrans digunakan untuk memproses pembayaran QRIS.
\item
  Algoritma Apriori digunakan untuk menganalisis pola pembelian pelanggan.
\end{itemize}

\begin{center}\rule{0.5\linewidth}{0.5pt}\end{center}

\subsection{2.3 Hak Akses Pengguna}\label{hak-akses-pengguna}

Sistem menerapkan Role-Based Access Control (RBAC) untuk membatasi akses pengguna sesuai perannya.

\subsubsection{Admin}\label{admin}

Admin memiliki hak akses penuh terhadap seluruh fitur sistem.

Hak akses:

\begin{itemize}
\tightlist
\item
  Dashboard
\item
  POS
\item
  Produk
\item
  Kategori
\item
  Pengguna
\item
  Laporan
\item
  Analytics
\end{itemize}

\subsubsection{Manager}\label{manager}

Manager memiliki akses untuk memantau operasional dan mengelola data produk.

Hak akses:

\begin{itemize}
\tightlist
\item
  Dashboard
\item
  POS
\item
  Produk
\item
  Kategori
\item
  Laporan
\item
  Analytics
\end{itemize}

\subsubsection{Cashier}\label{cashier}

Cashier bertugas melakukan transaksi penjualan.

Hak akses:

\begin{itemize}
\tightlist
\item
  Dashboard
\item
  POS
\end{itemize}

\begin{center}\rule{0.5\linewidth}{0.5pt}\end{center}

\subsection{2.4 Fitur Sistem}\label{fitur-sistem}

\subsubsection{2.4.1 Dashboard}\label{dashboard}

Dashboard berfungsi sebagai pusat informasi bisnis yang menampilkan ringkasan data penting.

Informasi yang ditampilkan meliputi:

\begin{itemize}
\tightlist
\item
  Pendapatan hari ini
\item
  Jumlah transaksi
\item
  Jumlah produk
\item
  Stok menipis
\item
  Transaksi terbaru
\end{itemize}

\begin{center}\rule{0.5\linewidth}{0.5pt}\end{center}

\subsubsection{2.4.2 Point of Sale (POS)}\label{point-of-sale-pos}

Modul POS digunakan oleh kasir untuk melakukan transaksi penjualan.

Fitur utama:

\begin{itemize}
\tightlist
\item
  Pencarian produk
\item
  Keranjang belanja
\item
  Diskon
\item
  Pajak
\item
  Checkout transaksi
\end{itemize}

\begin{center}\rule{0.5\linewidth}{0.5pt}\end{center}

\subsubsection{2.4.3 Pembayaran QRIS}\label{pembayaran-qris}

Sistem mendukung pembayaran QRIS melalui integrasi dengan Midtrans.

Alur pembayaran:

\begin{enumerate}
\def\labelenumi{\arabic{enumi}.}
\tightlist
\item
  Kasir memilih QRIS.
\item
  Sistem menghasilkan QR Code.
\item
  Pelanggan melakukan pembayaran.
\item
  Sistem memverifikasi pembayaran.
\item
  Transaksi dinyatakan berhasil.
\end{enumerate}

\begin{center}\rule{0.5\linewidth}{0.5pt}\end{center}

\subsubsection{2.4.4 Manajemen Produk}\label{manajemen-produk}

Modul produk digunakan untuk mengelola data barang yang dijual.

Fitur:

\begin{itemize}
\tightlist
\item
  Tambah produk
\item
  Edit produk
\item
  Hapus produk
\item
  Upload gambar
\item
  Pengelolaan stok
\end{itemize}

\begin{center}\rule{0.5\linewidth}{0.5pt}\end{center}

\subsubsection{2.4.5 Manajemen Kategori}\label{manajemen-kategori}

Kategori digunakan untuk mengelompokkan produk agar lebih mudah dikelola.

Fitur:

\begin{itemize}
\tightlist
\item
  Tambah kategori
\item
  Edit kategori
\item
  Hapus kategori
\end{itemize}

\begin{center}\rule{0.5\linewidth}{0.5pt}\end{center}

\subsubsection{2.4.6 Manajemen Pengguna}\label{manajemen-pengguna}

Modul pengguna digunakan untuk mengelola akun yang dapat mengakses sistem.

Fitur:

\begin{itemize}
\tightlist
\item
  Tambah pengguna
\item
  Edit pengguna
\item
  Hapus pengguna
\item
  Pengaturan role
\end{itemize}

\begin{center}\rule{0.5\linewidth}{0.5pt}\end{center}

\subsubsection{2.4.7 Laporan Penjualan}\label{laporan-penjualan}

Modul laporan digunakan untuk memantau performa bisnis.

Informasi yang tersedia:

\begin{itemize}
\tightlist
\item
  Total penjualan
\item
  Total transaksi
\item
  Produk terlaris
\item
  Grafik penjualan
\item
  Export CSV
\end{itemize}

\begin{center}\rule{0.5\linewidth}{0.5pt}\end{center}

\subsubsection{2.4.8 Market Basket Analysis}\label{market-basket-analysis}

Modul analytics menggunakan algoritma Apriori untuk menemukan hubungan antar produk yang sering dibeli secara bersamaan.

Output yang dihasilkan:

\begin{itemize}
\tightlist
\item
  Association Rules
\item
  Support
\item
  Confidence
\item
  Lift
\end{itemize}

Informasi tersebut dapat digunakan untuk strategi promosi dan bundling produk.

\begin{center}\rule{0.5\linewidth}{0.5pt}\end{center}

\subsection{2.5 Perancangan Database}\label{perancangan-database}

Database yang digunakan dalam sistem adalah MySQL.

\subsubsection{Tabel Users}\label{tabel-users}

Menyimpan data pengguna sistem.

{\def\LTcaptype{none} % do not increment counter
\begin{longtable}[]{@{}ll@{}}
\toprule\noalign{}
Field & Tipe \\
\midrule\noalign{}
\endhead
\bottomrule\noalign{}
\endlastfoot
id & bigint \\
name & varchar \\
email & varchar \\
password & varchar \\
role & varchar \\
\end{longtable}
}

\begin{center}\rule{0.5\linewidth}{0.5pt}\end{center}

\subsubsection{Tabel Categories}\label{tabel-categories}

Menyimpan data kategori produk.

{\def\LTcaptype{none} % do not increment counter
\begin{longtable}[]{@{}ll@{}}
\toprule\noalign{}
Field & Tipe \\
\midrule\noalign{}
\endhead
\bottomrule\noalign{}
\endlastfoot
id & bigint \\
name & varchar \\
description & text \\
\end{longtable}
}

\begin{center}\rule{0.5\linewidth}{0.5pt}\end{center}

\subsubsection{Tabel Products}\label{tabel-products}

Menyimpan data produk.

{\def\LTcaptype{none} % do not increment counter
\begin{longtable}[]{@{}ll@{}}
\toprule\noalign{}
Field & Tipe \\
\midrule\noalign{}
\endhead
\bottomrule\noalign{}
\endlastfoot
id & bigint \\
category\_id & bigint \\
name & varchar \\
price & decimal \\
stock & integer \\
\end{longtable}
}

\begin{center}\rule{0.5\linewidth}{0.5pt}\end{center}

\subsubsection{Tabel Orders}\label{tabel-orders}

Menyimpan transaksi penjualan.

{\def\LTcaptype{none} % do not increment counter
\begin{longtable}[]{@{}ll@{}}
\toprule\noalign{}
Field & Tipe \\
\midrule\noalign{}
\endhead
\bottomrule\noalign{}
\endlastfoot
id & bigint \\
user\_id & bigint \\
total\_price & decimal \\
payment\_method & varchar \\
payment\_status & varchar \\
\end{longtable}
}

\begin{center}\rule{0.5\linewidth}{0.5pt}\end{center}

\subsubsection{Tabel Order Items}\label{tabel-order-items}

Menyimpan detail item transaksi.

{\def\LTcaptype{none} % do not increment counter
\begin{longtable}[]{@{}ll@{}}
\toprule\noalign{}
Field & Tipe \\
\midrule\noalign{}
\endhead
\bottomrule\noalign{}
\endlastfoot
id & bigint \\
order\_id & bigint \\
product\_id & bigint \\
quantity & integer \\
price & decimal \\
\end{longtable}
}

\begin{center}\rule{0.5\linewidth}{0.5pt}\end{center}

\subsubsection{Tabel Inventory Batches}\label{tabel-inventory-batches}

Menyimpan data batch stok untuk implementasi FIFO.

{\def\LTcaptype{none} % do not increment counter
\begin{longtable}[]{@{}ll@{}}
\toprule\noalign{}
Field & Tipe \\
\midrule\noalign{}
\endhead
\bottomrule\noalign{}
\endlastfoot
id & bigint \\
product\_id & bigint \\
initial\_quantity & integer \\
remaining\_quantity & integer \\
buy\_price & decimal \\
\end{longtable}
}

\begin{center}\rule{0.5\linewidth}{0.5pt}\end{center}

\subsection{2.6 Entity Relationship Diagram (ERD)}\label{entity-relationship-diagram-erd}

Hubungan antar tabel dalam sistem dapat digambarkan sebagai berikut:

\begin{Shaded}
\begin{Highlighting}[]
\NormalTok{Users}
\NormalTok{  │}
\NormalTok{  └──── Orders}
\NormalTok{           │}
\NormalTok{           └──── Order Items}
\NormalTok{                     │}
\NormalTok{                     └──── Products}
\NormalTok{                               │}
\NormalTok{                               └──── Categories}

\NormalTok{Products}
\NormalTok{   │}
\NormalTok{   └──── Inventory Batches}
\end{Highlighting}
\end{Shaded}

\textbf{Tambahkan gambar ERD hasil draw.io atau StarUML pada bagian ini.}

\begin{center}\rule{0.5\linewidth}{0.5pt}\end{center}

\subsection{2.7 Use Case Diagram}\label{use-case-diagram}

Sistem melibatkan tiga aktor utama yaitu Admin, Manager, dan Cashier.

\subsubsection{Use Case Admin}\label{use-case-admin}

\begin{itemize}
\tightlist
\item
  Login
\item
  Kelola Produk
\item
  Kelola Kategori
\item
  Kelola Pengguna
\item
  Lihat Laporan
\item
  Lihat Analytics
\end{itemize}

\subsubsection{Use Case Manager}\label{use-case-manager}

\begin{itemize}
\tightlist
\item
  Login
\item
  Kelola Produk
\item
  Kelola Kategori
\item
  Lihat Laporan
\item
  Lihat Analytics
\end{itemize}

\subsubsection{Use Case Cashier}\label{use-case-cashier}

\begin{itemize}
\tightlist
\item
  Login
\item
  Melakukan Transaksi
\item
  Pembayaran Tunai
\item
  Pembayaran QRIS
\end{itemize}

\textbf{Tambahkan gambar Use Case Diagram pada bagian ini.}

\begin{center}\rule{0.5\linewidth}{0.5pt}\end{center}

\subsection{2.8 Ringkasan Perancangan Sistem}\label{ringkasan-perancangan-sistem}

Berdasarkan hasil perancangan yang telah dilakukan, sistem Brew \& Bytes POS dirancang sebagai aplikasi Point of Sale modern yang mampu mengelola transaksi penjualan, inventori, pembayaran digital, serta analisis pola pembelian pelanggan dalam satu platform terintegrasi.

Perancangan ini menjadi dasar dalam proses implementasi sistem yang akan dibahas pada bab berikutnya.

\section{BAB III}\label{bab-iii}

\section{IMPLEMENTASI SISTEM}\label{implementasi-sistem}

\subsection{3.1 Implementasi Sistem}\label{implementasi-sistem-1}

Tahap implementasi merupakan proses penerapan hasil perancangan sistem ke dalam bentuk aplikasi yang dapat digunakan oleh pengguna. Sistem Brew \& Bytes POS dikembangkan menggunakan Laravel sebagai backend, React sebagai frontend, MySQL sebagai basis data, serta Midtrans sebagai layanan pembayaran QRIS.

Implementasi dilakukan berdasarkan kebutuhan fungsional yang telah dirancang pada bab sebelumnya.

\begin{center}\rule{0.5\linewidth}{0.5pt}\end{center}

\subsection{3.2 Lingkungan Implementasi}\label{lingkungan-implementasi}

\subsubsection{3.2.1 Perangkat Keras}\label{perangkat-keras}

Perangkat keras yang digunakan selama proses pengembangan sistem adalah sebagai berikut:

{\def\LTcaptype{none} % do not increment counter
\begin{longtable}[]{@{}ll@{}}
\toprule\noalign{}
Komponen & Spesifikasi \\
\midrule\noalign{}
\endhead
\bottomrule\noalign{}
\endlastfoot
Processor & Intel Core i5 atau setara \\
RAM & 8 GB \\
Penyimpanan & SSD 256 GB \\
Sistem Operasi & Windows 11 \\
\end{longtable}
}

\subsubsection{3.2.2 Perangkat Lunak}\label{perangkat-lunak}

{\def\LTcaptype{none} % do not increment counter
\begin{longtable}[]{@{}ll@{}}
\toprule\noalign{}
Software & Versi \\
\midrule\noalign{}
\endhead
\bottomrule\noalign{}
\endlastfoot
PHP & 8.2 \\
Laravel & 12 \\
React & 19 \\
MySQL & 8 \\
Node.js & Terbaru \\
Visual Studio Code & Terbaru \\
Google Chrome & Terbaru \\
\end{longtable}
}

\begin{center}\rule{0.5\linewidth}{0.5pt}\end{center}

\section{3.3 Implementasi Antarmuka Sistem}\label{implementasi-antarmuka-sistem}

\subsection{3.3.1 Halaman Login}\label{halaman-login}

Halaman login digunakan sebagai gerbang autentikasi pengguna sebelum dapat mengakses sistem.

Fitur yang tersedia pada halaman login meliputi:

\begin{itemize}
\tightlist
\item
  Input email
\item
  Input password
\item
  Validasi login
\item
  Hak akses berdasarkan role pengguna
\end{itemize}

\subsubsection{Tampilan Halaman Login}\label{tampilan-halaman-login}

\textbf{Gambar 3.1 Halaman Login}

\emph{(Tambahkan screenshot halaman login di sini)}

\subsubsection{Hasil Implementasi}\label{hasil-implementasi}

Sistem berhasil melakukan autentikasi pengguna berdasarkan data yang tersimpan pada database dan mengarahkan pengguna ke dashboard sesuai hak akses yang dimiliki.

\begin{center}\rule{0.5\linewidth}{0.5pt}\end{center}

\subsection{3.3.2 Dashboard}\label{dashboard-1}

Dashboard merupakan halaman utama yang menampilkan ringkasan kondisi bisnis secara real-time.

Informasi yang ditampilkan meliputi:

\begin{itemize}
\tightlist
\item
  Total pendapatan hari ini
\item
  Total transaksi
\item
  Jumlah produk
\item
  Stok menipis
\item
  Transaksi terbaru
\end{itemize}

\subsubsection{Tampilan Dashboard}\label{tampilan-dashboard}

\textbf{Gambar 3.2 Dashboard}

\emph{(Tambahkan screenshot dashboard di sini)}

\subsubsection{Hasil Implementasi}\label{hasil-implementasi-1}

Dashboard berhasil menyajikan informasi penting yang dibutuhkan pengguna untuk memantau kondisi operasional bisnis.

\begin{center}\rule{0.5\linewidth}{0.5pt}\end{center}

\subsection{3.3.3 Point of Sale (POS)}\label{point-of-sale-pos-1}

Modul Point of Sale digunakan untuk melakukan transaksi penjualan.

Fitur utama:

\begin{itemize}
\tightlist
\item
  Katalog produk
\item
  Pencarian produk
\item
  Keranjang belanja
\item
  Diskon
\item
  Pajak
\item
  Checkout transaksi
\end{itemize}

\subsubsection{Tampilan POS}\label{tampilan-pos}

\textbf{Gambar 3.3 Halaman POS}

\emph{(Tambahkan screenshot POS di sini)}

\subsubsection{Hasil Implementasi}\label{hasil-implementasi-2}

Kasir dapat melakukan transaksi penjualan secara cepat dan efisien melalui antarmuka yang sederhana dan mudah digunakan.

\begin{center}\rule{0.5\linewidth}{0.5pt}\end{center}

\subsection{3.3.4 Keranjang Belanja}\label{keranjang-belanja}

Keranjang belanja digunakan untuk menampilkan daftar produk yang dipilih pelanggan sebelum transaksi diselesaikan.

Informasi yang ditampilkan:

\begin{itemize}
\tightlist
\item
  Nama produk
\item
  Harga produk
\item
  Jumlah produk
\item
  Subtotal
\item
  Diskon
\item
  Pajak
\item
  Total pembayaran
\end{itemize}

\subsubsection{Tampilan Keranjang}\label{tampilan-keranjang}

\textbf{Gambar 3.4 Keranjang Belanja}

\emph{(Tambahkan screenshot keranjang belanja di sini)}

\begin{center}\rule{0.5\linewidth}{0.5pt}\end{center}

\subsection{3.3.5 Pembayaran QRIS}\label{pembayaran-qris-1}

Sistem mendukung pembayaran digital menggunakan QRIS melalui integrasi dengan Midtrans.

Alur pembayaran:

\begin{enumerate}
\def\labelenumi{\arabic{enumi}.}
\tightlist
\item
  Kasir memilih metode pembayaran QRIS.
\item
  Sistem menghasilkan QR Code.
\item
  Pelanggan melakukan pembayaran.
\item
  Sistem melakukan verifikasi transaksi.
\item
  Status transaksi berubah menjadi berhasil.
\end{enumerate}

\subsubsection{Tampilan QRIS}\label{tampilan-qris}

\textbf{Gambar 3.5 Pembayaran QRIS}

\emph{(Tambahkan screenshot QRIS di sini)}

\subsubsection{Hasil Implementasi}\label{hasil-implementasi-3}

Sistem berhasil menghasilkan QR Code pembayaran dan menerima konfirmasi pembayaran secara otomatis.

\begin{center}\rule{0.5\linewidth}{0.5pt}\end{center}

\subsection{3.3.6 Manajemen Produk}\label{manajemen-produk-1}

Halaman produk digunakan untuk mengelola data produk yang dijual.

Fitur:

\begin{itemize}
\tightlist
\item
  Tambah produk
\item
  Edit produk
\item
  Hapus produk
\item
  Upload gambar produk
\item
  Pengelolaan stok
\end{itemize}

\subsubsection{Tampilan Produk}\label{tampilan-produk}

\textbf{Gambar 3.6 Manajemen Produk}

\emph{(Tambahkan screenshot halaman produk di sini)}

\subsubsection{Hasil Implementasi}\label{hasil-implementasi-4}

Data produk dapat dikelola secara terpusat dan langsung digunakan pada proses transaksi.

\begin{center}\rule{0.5\linewidth}{0.5pt}\end{center}

\subsection{3.3.7 Manajemen Kategori}\label{manajemen-kategori-1}

Halaman kategori digunakan untuk mengelompokkan produk berdasarkan jenisnya.

Fitur:

\begin{itemize}
\tightlist
\item
  Tambah kategori
\item
  Edit kategori
\item
  Hapus kategori
\end{itemize}

\subsubsection{Tampilan Kategori}\label{tampilan-kategori}

\textbf{Gambar 3.7 Manajemen Kategori}

\emph{(Tambahkan screenshot kategori di sini)}

\begin{center}\rule{0.5\linewidth}{0.5pt}\end{center}

\subsection{3.3.8 Manajemen Pengguna}\label{manajemen-pengguna-1}

Modul pengguna digunakan untuk mengelola akun yang dapat mengakses sistem.

Hak akses yang tersedia:

\begin{itemize}
\tightlist
\item
  Admin
\item
  Manager
\item
  Cashier
\end{itemize}

\subsubsection{Tampilan Manajemen Pengguna}\label{tampilan-manajemen-pengguna}

\textbf{Gambar 3.8 Manajemen Pengguna}

\emph{(Tambahkan screenshot user management di sini)}

\subsubsection{Hasil Implementasi}\label{hasil-implementasi-5}

Admin dapat mengelola akun pengguna dan menentukan hak akses sesuai kebutuhan operasional.

\begin{center}\rule{0.5\linewidth}{0.5pt}\end{center}

\subsection{3.3.9 Laporan Penjualan}\label{laporan-penjualan-1}

Halaman laporan digunakan untuk memantau performa bisnis berdasarkan data transaksi.

Informasi yang tersedia:

\begin{itemize}
\tightlist
\item
  Total penjualan
\item
  Total transaksi
\item
  Produk terlaris
\item
  Grafik penjualan
\end{itemize}

\subsubsection{Tampilan Laporan}\label{tampilan-laporan}

\textbf{Gambar 3.9 Laporan Penjualan}

\emph{(Tambahkan screenshot laporan di sini)}

\subsubsection{Hasil Implementasi}\label{hasil-implementasi-6}

Laporan berhasil menampilkan informasi penjualan secara visual sehingga memudahkan proses analisis bisnis.

\begin{center}\rule{0.5\linewidth}{0.5pt}\end{center}

\subsection{3.3.10 Analytics Apriori}\label{analytics-apriori}

Halaman Analytics digunakan untuk menampilkan hasil analisis pola pembelian pelanggan menggunakan algoritma Apriori.

Informasi yang ditampilkan:

\begin{itemize}
\tightlist
\item
  Association Rules
\item
  Support
\item
  Confidence
\item
  Lift
\end{itemize}

\subsubsection{Tampilan Analytics}\label{tampilan-analytics}

\textbf{Gambar 3.10 Analytics Apriori}

\emph{(Tambahkan screenshot analytics di sini)}

\subsubsection{Hasil Implementasi}\label{hasil-implementasi-7}

Sistem berhasil menghasilkan aturan asosiasi berdasarkan data transaksi yang tersimpan.

\begin{center}\rule{0.5\linewidth}{0.5pt}\end{center}

\section{3.4 Implementasi FIFO}\label{implementasi-fifo}

Sistem menerapkan metode FIFO (First In First Out) dalam pengelolaan stok.

Metode FIFO memastikan stok yang pertama kali masuk akan menjadi stok yang pertama kali digunakan.

\subsubsection{Contoh Implementasi}\label{contoh-implementasi}

{\def\LTcaptype{none} % do not increment counter
\begin{longtable}[]{@{}ll@{}}
\toprule\noalign{}
Batch & Jumlah \\
\midrule\noalign{}
\endhead
\bottomrule\noalign{}
\endlastfoot
Batch 1 & 50 \\
Batch 2 & 50 \\
\end{longtable}
}

Jika terjadi penjualan sebanyak 60 unit maka:

{\def\LTcaptype{none} % do not increment counter
\begin{longtable}[]{@{}ll@{}}
\toprule\noalign{}
Batch & Sisa \\
\midrule\noalign{}
\endhead
\bottomrule\noalign{}
\endlastfoot
Batch 1 & 0 \\
Batch 2 & 40 \\
\end{longtable}
}

\subsubsection{Hasil Implementasi}\label{hasil-implementasi-8}

Metode FIFO berhasil diterapkan untuk menjaga akurasi pengelolaan persediaan.

\begin{center}\rule{0.5\linewidth}{0.5pt}\end{center}

\section{3.5 Implementasi Market Basket Analysis}\label{implementasi-market-basket-analysis}

Market Basket Analysis diterapkan menggunakan algoritma Apriori.

Tujuannya adalah menemukan hubungan antar produk yang sering dibeli secara bersamaan.

\subsubsection{Contoh Data Transaksi}\label{contoh-data-transaksi}

{\def\LTcaptype{none} % do not increment counter
\begin{longtable}[]{@{}ll@{}}
\toprule\noalign{}
Transaksi & Produk \\
\midrule\noalign{}
\endhead
\bottomrule\noalign{}
\endlastfoot
T1 & Espresso, Croissant \\
T2 & Espresso, Brownies \\
T3 & Espresso, Croissant \\
T4 & Latte, Brownies \\
T5 & Espresso, Croissant \\
\end{longtable}
}

\subsubsection{Contoh Hasil Analisis}\label{contoh-hasil-analisis}

{\def\LTcaptype{none} % do not increment counter
\begin{longtable}[]{@{}lll@{}}
\toprule\noalign{}
Rule & Support & Confidence \\
\midrule\noalign{}
\endhead
\bottomrule\noalign{}
\endlastfoot
Espresso → Croissant & 60\% & 75\% \\
Croissant → Espresso & 60\% & 100\% \\
\end{longtable}
}

\subsubsection{Hasil Implementasi}\label{hasil-implementasi-9}

Hasil analisis dapat digunakan sebagai dasar dalam penyusunan strategi promosi dan paket bundling produk.

\begin{center}\rule{0.5\linewidth}{0.5pt}\end{center}

\section{3.6 Ringkasan Implementasi}\label{ringkasan-implementasi}

Berdasarkan hasil implementasi yang telah dilakukan, seluruh modul utama sistem berhasil dibangun dan dapat berjalan sesuai kebutuhan yang telah ditentukan.

Modul yang berhasil diimplementasikan meliputi:

\begin{itemize}
\tightlist
\item
  Login dan Manajemen Pengguna
\item
  Dashboard
\item
  Point of Sale (POS)
\item
  Manajemen Produk
\item
  Manajemen Kategori
\item
  Pembayaran QRIS
\item
  Laporan Penjualan
\item
  FIFO Inventory
\item
  Analytics Apriori
\item
  Sinkronisasi Offline
\end{itemize}

Dengan demikian sistem Brew \& Bytes POS telah berhasil dikembangkan sebagai aplikasi Point of Sale berbasis web yang mampu mendukung operasional bisnis secara terintegrasi.

\section{BAB IV}\label{bab-iv}

\section{PENGUJIAN DAN KESIMPULAN}\label{pengujian-dan-kesimpulan}

\subsection{4.1 Pengujian Sistem}\label{pengujian-sistem}

Pengujian sistem dilakukan untuk memastikan seluruh fitur yang terdapat pada Brew \& Bytes POS dapat berjalan sesuai dengan kebutuhan yang telah dirancang. Metode pengujian yang digunakan adalah \textbf{Black Box Testing}, yaitu metode pengujian yang berfokus pada fungsi sistem tanpa melihat implementasi kode program.

Tujuan pengujian adalah untuk memastikan bahwa setiap fitur dapat menerima input, memproses data, dan menghasilkan output yang sesuai dengan kebutuhan pengguna.

\begin{center}\rule{0.5\linewidth}{0.5pt}\end{center}

\section{4.2 Pengujian Fitur Login}\label{pengujian-fitur-login}

Pengujian dilakukan untuk memastikan sistem autentikasi berjalan dengan baik.

{\def\LTcaptype{none} % do not increment counter
\begin{longtable}[]{@{}
  >{\raggedright\arraybackslash}p{(\linewidth - 6\tabcolsep) * \real{0.0270}}
  >{\raggedright\arraybackslash}p{(\linewidth - 6\tabcolsep) * \real{0.5000}}
  >{\raggedright\arraybackslash}p{(\linewidth - 6\tabcolsep) * \real{0.3649}}
  >{\raggedright\arraybackslash}p{(\linewidth - 6\tabcolsep) * \real{0.1081}}@{}}
\toprule\noalign{}
\begin{minipage}[b]{\linewidth}\raggedright
No
\end{minipage} & \begin{minipage}[b]{\linewidth}\raggedright
Skenario Pengujian
\end{minipage} & \begin{minipage}[b]{\linewidth}\raggedright
Hasil Yang Diharapkan
\end{minipage} & \begin{minipage}[b]{\linewidth}\raggedright
Hasil
\end{minipage} \\
\midrule\noalign{}
\endhead
\bottomrule\noalign{}
\endlastfoot
1 & Login dengan email dan password valid & Berhasil masuk ke dashboard & Berhasil \\
2 & Login dengan password salah & Menampilkan pesan error & Berhasil \\
3 & Login dengan email kosong & Menampilkan validasi & Berhasil \\
4 & Login dengan password kosong & Menampilkan validasi & Berhasil \\
\end{longtable}
}

\subsubsection{Hasil Pengujian}\label{hasil-pengujian}

Berdasarkan pengujian yang dilakukan, fitur login berhasil berjalan sesuai dengan kebutuhan sistem.

\begin{center}\rule{0.5\linewidth}{0.5pt}\end{center}

\section{4.3 Pengujian Manajemen Produk}\label{pengujian-manajemen-produk}

Pengujian dilakukan pada fitur pengelolaan produk.

{\def\LTcaptype{none} % do not increment counter
\begin{longtable}[]{@{}llll@{}}
\toprule\noalign{}
No & Skenario Pengujian & Hasil Yang Diharapkan & Hasil \\
\midrule\noalign{}
\endhead
\bottomrule\noalign{}
\endlastfoot
1 & Menambah produk baru & Produk tersimpan & Berhasil \\
2 & Mengubah data produk & Data diperbarui & Berhasil \\
3 & Menghapus produk & Produk terhapus & Berhasil \\
4 & Upload gambar produk & Gambar tersimpan & Berhasil \\
\end{longtable}
}

\subsubsection{Hasil Pengujian}\label{hasil-pengujian-1}

Seluruh fitur manajemen produk berhasil berjalan dengan baik.

\begin{center}\rule{0.5\linewidth}{0.5pt}\end{center}

\section{4.4 Pengujian Manajemen Kategori}\label{pengujian-manajemen-kategori}

{\def\LTcaptype{none} % do not increment counter
\begin{longtable}[]{@{}llll@{}}
\toprule\noalign{}
No & Skenario Pengujian & Hasil Yang Diharapkan & Hasil \\
\midrule\noalign{}
\endhead
\bottomrule\noalign{}
\endlastfoot
1 & Menambah kategori & Data tersimpan & Berhasil \\
2 & Mengubah kategori & Data diperbarui & Berhasil \\
3 & Menghapus kategori & Data terhapus & Berhasil \\
\end{longtable}
}

\subsubsection{Hasil Pengujian}\label{hasil-pengujian-2}

Fitur kategori berhasil dijalankan sesuai dengan kebutuhan.

\begin{center}\rule{0.5\linewidth}{0.5pt}\end{center}

\section{4.5 Pengujian Transaksi Point of Sale}\label{pengujian-transaksi-point-of-sale}

Pengujian dilakukan terhadap proses transaksi penjualan.

{\def\LTcaptype{none} % do not increment counter
\begin{longtable}[]{@{}
  >{\raggedright\arraybackslash}p{(\linewidth - 6\tabcolsep) * \real{0.0317}}
  >{\raggedright\arraybackslash}p{(\linewidth - 6\tabcolsep) * \real{0.4921}}
  >{\raggedright\arraybackslash}p{(\linewidth - 6\tabcolsep) * \real{0.3492}}
  >{\raggedright\arraybackslash}p{(\linewidth - 6\tabcolsep) * \real{0.1270}}@{}}
\toprule\noalign{}
\begin{minipage}[b]{\linewidth}\raggedright
No
\end{minipage} & \begin{minipage}[b]{\linewidth}\raggedright
Skenario Pengujian
\end{minipage} & \begin{minipage}[b]{\linewidth}\raggedright
Hasil Yang Diharapkan
\end{minipage} & \begin{minipage}[b]{\linewidth}\raggedright
Hasil
\end{minipage} \\
\midrule\noalign{}
\endhead
\bottomrule\noalign{}
\endlastfoot
1 & Menambahkan produk ke keranjang & Produk masuk keranjang & Berhasil \\
2 & Mengubah jumlah produk & Total diperbarui & Berhasil \\
3 & Menghapus produk dari keranjang & Produk terhapus & Berhasil \\
4 & Checkout transaksi & Transaksi tersimpan & Berhasil \\
\end{longtable}
}

\subsubsection{Hasil Pengujian}\label{hasil-pengujian-3}

Modul Point of Sale berhasil memproses transaksi penjualan dengan baik.

\begin{center}\rule{0.5\linewidth}{0.5pt}\end{center}

\section{4.6 Pengujian Pembayaran QRIS}\label{pengujian-pembayaran-qris}

Pengujian dilakukan pada integrasi Payment Gateway Midtrans.

{\def\LTcaptype{none} % do not increment counter
\begin{longtable}[]{@{}
  >{\raggedright\arraybackslash}p{(\linewidth - 6\tabcolsep) * \real{0.0278}}
  >{\raggedright\arraybackslash}p{(\linewidth - 6\tabcolsep) * \real{0.3056}}
  >{\raggedright\arraybackslash}p{(\linewidth - 6\tabcolsep) * \real{0.5556}}
  >{\raggedright\arraybackslash}p{(\linewidth - 6\tabcolsep) * \real{0.1111}}@{}}
\toprule\noalign{}
\begin{minipage}[b]{\linewidth}\raggedright
No
\end{minipage} & \begin{minipage}[b]{\linewidth}\raggedright
Skenario Pengujian
\end{minipage} & \begin{minipage}[b]{\linewidth}\raggedright
Hasil Yang Diharapkan
\end{minipage} & \begin{minipage}[b]{\linewidth}\raggedright
Hasil
\end{minipage} \\
\midrule\noalign{}
\endhead
\bottomrule\noalign{}
\endlastfoot
1 & Generate QRIS & QR Code berhasil dibuat & Berhasil \\
2 & Pembayaran berhasil & Status transaksi menjadi Paid & Berhasil \\
3 & Pembayaran dibatalkan & Status transaksi menjadi Failed & Berhasil \\
4 & Pembayaran kedaluwarsa & Status transaksi berubah menjadi Expired & Berhasil \\
\end{longtable}
}

\subsubsection{Hasil Pengujian}\label{hasil-pengujian-4}

Integrasi QRIS berhasil berjalan sesuai dengan alur pembayaran yang dirancang.

\begin{center}\rule{0.5\linewidth}{0.5pt}\end{center}

\section{4.7 Pengujian FIFO}\label{pengujian-fifo}

Pengujian dilakukan untuk memastikan sistem mengurangi stok berdasarkan urutan batch yang masuk terlebih dahulu.

\subsubsection{Data Awal}\label{data-awal}

{\def\LTcaptype{none} % do not increment counter
\begin{longtable}[]{@{}ll@{}}
\toprule\noalign{}
Batch & Jumlah \\
\midrule\noalign{}
\endhead
\bottomrule\noalign{}
\endlastfoot
Batch 1 & 50 \\
Batch 2 & 50 \\
\end{longtable}
}

Total stok: 100 unit

\subsubsection{Simulasi Penjualan}\label{simulasi-penjualan}

Penjualan sebanyak 60 unit.

\subsubsection{Hasil}\label{hasil}

{\def\LTcaptype{none} % do not increment counter
\begin{longtable}[]{@{}ll@{}}
\toprule\noalign{}
Batch & Sisa \\
\midrule\noalign{}
\endhead
\bottomrule\noalign{}
\endlastfoot
Batch 1 & 0 \\
Batch 2 & 40 \\
\end{longtable}
}

\subsubsection{Analisis}\label{analisis}

Hasil menunjukkan bahwa sistem berhasil menerapkan metode FIFO dengan benar. Batch pertama digunakan terlebih dahulu sebelum mengambil stok dari batch berikutnya.

\begin{center}\rule{0.5\linewidth}{0.5pt}\end{center}

\section{4.8 Pengujian Analytics Apriori}\label{pengujian-analytics-apriori}

Pengujian dilakukan untuk memastikan sistem dapat menghasilkan aturan asosiasi berdasarkan data transaksi.

\subsubsection{Data Sampel}\label{data-sampel}

{\def\LTcaptype{none} % do not increment counter
\begin{longtable}[]{@{}ll@{}}
\toprule\noalign{}
Transaksi & Produk \\
\midrule\noalign{}
\endhead
\bottomrule\noalign{}
\endlastfoot
T1 & Espresso, Croissant \\
T2 & Espresso, Brownies \\
T3 & Espresso, Croissant \\
T4 & Latte, Brownies \\
T5 & Espresso, Croissant \\
\end{longtable}
}

\subsubsection{Hasil Analisis}\label{hasil-analisis}

{\def\LTcaptype{none} % do not increment counter
\begin{longtable}[]{@{}lll@{}}
\toprule\noalign{}
Rule & Support & Confidence \\
\midrule\noalign{}
\endhead
\bottomrule\noalign{}
\endlastfoot
Espresso → Croissant & 60\% & 75\% \\
Croissant → Espresso & 60\% & 100\% \\
\end{longtable}
}

\subsubsection{Analisis}\label{analisis-1}

Berdasarkan hasil pengujian, algoritma Apriori berhasil menemukan hubungan antar produk yang sering dibeli secara bersamaan.

Informasi tersebut dapat digunakan untuk:

\begin{itemize}
\tightlist
\item
  Strategi bundling produk
\item
  Cross-selling
\item
  Penempatan produk
\item
  Program promosi
\end{itemize}

\begin{center}\rule{0.5\linewidth}{0.5pt}\end{center}

\section{4.9 Pengujian Mode Offline}\label{pengujian-mode-offline}

Pengujian dilakukan untuk memastikan sistem tetap dapat digunakan ketika koneksi internet terputus.

{\def\LTcaptype{none} % do not increment counter
\begin{longtable}[]{@{}
  >{\raggedright\arraybackslash}p{(\linewidth - 6\tabcolsep) * \real{0.0312}}
  >{\raggedright\arraybackslash}p{(\linewidth - 6\tabcolsep) * \real{0.3750}}
  >{\raggedright\arraybackslash}p{(\linewidth - 6\tabcolsep) * \real{0.4688}}
  >{\raggedright\arraybackslash}p{(\linewidth - 6\tabcolsep) * \real{0.1250}}@{}}
\toprule\noalign{}
\begin{minipage}[b]{\linewidth}\raggedright
No
\end{minipage} & \begin{minipage}[b]{\linewidth}\raggedright
Skenario Pengujian
\end{minipage} & \begin{minipage}[b]{\linewidth}\raggedright
Hasil Yang Diharapkan
\end{minipage} & \begin{minipage}[b]{\linewidth}\raggedright
Hasil
\end{minipage} \\
\midrule\noalign{}
\endhead
\bottomrule\noalign{}
\endlastfoot
1 & Internet terputus & Sistem tetap berjalan & Berhasil \\
2 & Simpan transaksi offline & Data tersimpan lokal & Berhasil \\
3 & Internet kembali aktif & Sinkronisasi otomatis berjalan & Berhasil \\
\end{longtable}
}

\subsubsection{Hasil Pengujian}\label{hasil-pengujian-5}

Fitur offline berhasil berjalan sesuai perancangan dan mampu menyimpan transaksi sementara sebelum dilakukan sinkronisasi ke server.

\begin{center}\rule{0.5\linewidth}{0.5pt}\end{center}

\section{4.10 Ringkasan Hasil Pengujian}\label{ringkasan-hasil-pengujian}

Berdasarkan seluruh pengujian yang telah dilakukan, seluruh fitur utama sistem berhasil berjalan sesuai kebutuhan.

Ringkasan hasil pengujian dapat dilihat pada tabel berikut.

{\def\LTcaptype{none} % do not increment counter
\begin{longtable}[]{@{}ll@{}}
\toprule\noalign{}
Modul & Status \\
\midrule\noalign{}
\endhead
\bottomrule\noalign{}
\endlastfoot
Login & Berhasil \\
Dashboard & Berhasil \\
POS & Berhasil \\
Produk & Berhasil \\
Kategori & Berhasil \\
Pengguna & Berhasil \\
QRIS & Berhasil \\
FIFO & Berhasil \\
Analytics Apriori & Berhasil \\
Offline Sync & Berhasil \\
\end{longtable}
}

\begin{center}\rule{0.5\linewidth}{0.5pt}\end{center}

\section{4.11 Kesimpulan}\label{kesimpulan}

Berdasarkan hasil pengembangan dan pengujian yang telah dilakukan, dapat disimpulkan bahwa proyek Brew \& Bytes POS berhasil dibangun sesuai dengan kebutuhan yang telah ditentukan.

Sistem mampu mendukung proses transaksi penjualan, pengelolaan produk, pengelolaan stok, pembayaran digital melalui QRIS, penyusunan laporan penjualan, serta analisis pola pembelian pelanggan menggunakan algoritma Apriori.

Penerapan metode FIFO berhasil membantu pengelolaan persediaan barang, sedangkan implementasi algoritma Apriori mampu menghasilkan informasi yang dapat digunakan untuk mendukung strategi pemasaran dan pengambilan keputusan bisnis.

Secara keseluruhan, Brew \& Bytes POS berhasil menjadi solusi terintegrasi yang dapat membantu operasional kafe dan restoran secara lebih efektif, efisien, dan berbasis data.
